%%%%%%%%%%%%%%%%
%%% deut 29 %%%
%%%%%%%%%%%%%%%%

\Chapter{29}

\Verse{1}
{
	\hDtr1{וַיִּקְרָא מֹשֶׁה אֶל כּל יִשְׂרָאֵל וַיֹּאמֶר אֲלֵהֶם אַתֶּם רְאִיתֶם אֵת כּל אֲשֶׁר עָשָׂה יְהֹוָה לְעֵינֵיכֶם בְּאֶרֶץ מִצְרַיִם לְפַרְעֹה וּלְכל עֲבָדָיו וּלְכל אַרְצוֹ׃}
}
{
	\eDtr1{
		29 And Moses called all of Israel and said to them,
		“You’ve seen everything that YHWH did before your eyes in
		the land of Egypt to Pharaoh and to all his servants and
		to all his land,
	}
}
{
}

\Verse{2}
{
	\hDtr1{הַמַּסּוֹת הַגְּדֹלֹת אֲשֶׁר רָאוּ עֵינֶיךָ הָאֹתֹת וְהַמֹּפְתִים הַגְּדֹלִים הָהֵם׃}
}
{
	\eDtr1{
		the great tests that your eyes saw, those great signs and
		wonders.
	}
}
{
}

\Verse{3}
{
	\hDtr1{וְלֹא נָתַן יְהֹוָה לָכֶם לֵב לָדַעַת וְעֵינַיִם לִרְאוֹת וְאזְנַיִם לִשְׁמֹעַ עַד הַיּוֹם הַזֶּה׃}
}
{
	\eDtr1{
		But YHWH did not give you a heart to know and eyes to see
		and ears to hear until this day.
	}
}
{
%	\footnote{
%		did not give you … until this day. This is an old enigma: have they not
%		been able to understand the meaning of all the great wonders they have
%		witnessed until now? I understand this to mean just that. Moses has just
%		finished the long speech that takes up nearly all of Deuteronomy. Now he
%		is beginning a much shorter speech, and he is telling them that he is
%		about to reveal to them, in his last message to them, the key to
%		understanding what their experience means. What follows is a speech like
%		no other he has ever made in the Torah. Until now he has focused on
%		facts: points of history, specific laws. But, starting in a few verses,
%		he will turn philosophical, moral, and spiritual—and his message is
%		formulated in some of the most beautiful words in the Torah.
%	}
}

\Verse{4}
{
	\hDtr1{וָאוֹלֵךְ אֶתְכֶם אַרְבָּעִים שָׁנָה בַּמִּדְבָּר לֹא בָלוּ שַׂלְמֹתֵיכֶם מֵעֲלֵיכֶם וְנַעַלְךָ לֹא בָלְתָה מֵעַל רַגְלֶךָ׃}
}
{
	\eDtr1{
		As I led you forty years in the wilderness, your clothing
		did not become worn on you, and your shoe did not become
		worn on your foot,
	}
}
{
}

\Verse{5}
{
	\hDtr1{לֶחֶם לֹא אֲכַלְתֶּם וְיַיִן וְשֵׁכָר לֹא שְׁתִיתֶם לְמַעַן תֵּדְעוּ כִּי אֲנִי יְהֹוָה אֱלֹהֵיכֶם׃}
}
{
	\eDtr1{
		you didn’t eat bread, and you didn’t drink wine and beer,
		so you would know that ‘I am YHWH your {\God}.’
	}
}
{
}

\Verse{6}
{
	\hDtr1{וַתָּבֹאוּ אֶל הַמָּקוֹם הַזֶּה וַיֵּצֵא סִיחֹן מֶלֶךְ חֶשְׁבּוֹן וְעוֹג מֶלֶךְ הַבָּשָׁן לִקְרָאתֵנוּ לַמִּלְחָמָה וַנַּכֵּם׃}
}
{
	\eDtr1{
		And you came to this place. (And Sihon, king of Heshbon,
		and Og, king of Bashan, came out at us for war, and we
		struck them.
	}
}
{
}

\Verse{7}
{
	\hDtr1{וַנִּקַּח אֶת אַרְצָם וַנִּתְּנָהּ לְנַחֲלָה לָראוּבֵנִי וְלַגָּדִי וְלַחֲצִי שֵׁבֶט הַמְנַשִּׁי׃}
}
{
	\eDtr1{
		And we took their land and gave it as a legacy to the
		Reubenite and to the Gadite and to half of the Manassite
		tribe.)
	}
}
{
}

\Verse{8}
{
	\hDtr1{וּשְׁמַרְתֶּם אֶת דִּבְרֵי הַבְּרִית הַזֹּאת וַעֲשִׂיתֶם אֹתָם לְמַעַן תַּשְׂכִּילוּ אֵת כּל אֲשֶׁר תַּעֲשׂוּן׃}
}
{
	\eDtr1{
		And you shall observe the words of this covenant and do
		them so that you’ll understand all that you’ll do.
	}
}
{
%	\footnote{
%		understand. Yet again the words of the beginning of the Torah return in
%		the Torah’s closing portions. In the garden of Eden, the woman is
%		attracted to the fruit of the tree of knowledge of good and bad because
%		it will “bring about understanding” (Hebrew ha[image]kîl; 3:6). Humans
%		acquire that ability—understanding—when they eat the forbidden fruit.
%		The word then does not occur again in the Torah until here. Now Moses
%		tells them that they can use the divine power of understanding for a
%		purpose: comprehending the meaning of what they have experienced and of
%		what they will do. The past and the future, the miracles and the
%		performance of the covenant’s commandments, are all within humans’ power
%		to grasp. What began as something taken in violation of a divine command
%		now becomes something positive. Humans are not limited to following the
%		commandments without comprehension. They can use what they took from the
%		divine realm for their own good and for their relationship with their
%		God.
%	}
}

\Verse{9}
{
	\hDtr1{אַתֶּם נִצָּבִים הַיּוֹם כֻּלְּכֶם לִפְנֵי יְהֹוָה אֱלֹהֵיכֶם רָאשֵׁיכֶם שִׁבְטֵיכֶם זִקְנֵיכֶם וְשֹׁטְרֵיכֶם כֹּל אִישׁ יִשְׂרָאֵל׃}
}
{
	\eDtr1{
		STANDING “You’re standing today, all of you, in front of
		YHWH, your {\God}—your heads, your tribes, your elders, and
		your officers, every man of Israel,
	}
}
{
}

\Verse{10}
{
	\hDtr1{טַפְּכֶם נְשֵׁיכֶם וְגֵרְךָ אֲשֶׁר בְּקֶרֶב מַחֲנֶיךָ מֵחֹטֵב עֵצֶיךָ עַד שֹׁאֵב מֵימֶיךָ׃}
}
{
	\eDtr1{
		your infants, your women, and your alien who is in your
		camps, from one who cuts your wood to one who draws your
		water—
	}
}
{
}

\Verse{11}
{
	\hDtr1{לְעבְרְךָ בִּבְרִית יְהֹוָה אֱלֹהֶיךָ וּבְאָלָתוֹ אֲשֶׁר יְהֹוָה אֱלֹהֶיךָ כֹּרֵת עִמְּךָ הַיּוֹם׃}
}
{
	\eDtr1{
		for you to enter into the covenant of YHWH, your {\God}, and
		into His oath, which YHWH, your {\God}, is making with you
		today,
	}
}
{
}

\Verse{12}
{
	\hDtr1{לְמַעַן הָקִים אֹתְךָ הַיּוֹם לוֹ לְעָם וְהוּא יִהְיֶה לְּךָ לֵאלֹהִים כַּאֲשֶׁר דִּבֶּר לָךְ וְכַאֲשֶׁר נִשְׁבַּע לַאֲבֹתֶיךָ לְאַבְרָהָם לְיִצְחָק וּלְיַעֲקֹב׃}
}
{
	\eDtr1{
		in order to establish you today for Him as a people, and
		He will be a {\God} to you, as He spoke to you and as He
		swore to your fathers: to Abraham, to Isaac, and to Jacob.
	}
}
{
}

\Verse{13}
{
	\hDtr1{וְלֹא אִתְּכֶם לְבַדְּכֶם אָנֹכִי כֹּרֵת אֶת הַבְּרִית הַזֹּאת וְאֶת הָאָלָה הַזֹּאת׃}
}
{
	\eDtr1{
		And I am not making this covenant and this oath with you
		alone,
	}
}
{
}

\Verse{14}
{
	\hDtr1{כִּי אֶת אֲשֶׁר יֶשְׁנוֹ פֹּה עִמָּנוּ עֹמֵד הַיּוֹם לִפְנֵי יְהֹוָה אֱלֹהֵינוּ וְאֵת אֲשֶׁר אֵינֶנּוּ פֹּה עִמָּנוּ הַיּוֹם׃}
}
{
	\eDtr1{
		but with the one who is here standing with us today in
		front of YHWH, our {\God}, and with the one who isn’t here
		with us today.
	}
}
{
}

\Verse{15}
{
	\hDtr1{כִּי אַתֶּם יְדַעְתֶּם אֵת אֲשֶׁר יָשַׁבְנוּ בְּאֶרֶץ מִצְרָיִם וְאֵת אֲשֶׁר עָבַרְנוּ בְּקֶרֶב הַגּוֹיִם אֲשֶׁר עֲבַרְתֶּם׃}
}
{
	\eDtr1{
		Because you know that we lived in the land of Egypt and
		that we passed among the nations that you passed.
	}
}
{
}

\Verse{16}
{
	\hDtr1{וַתִּרְאוּ אֶת שִׁקּוּצֵיהֶם וְאֵת גִּלֻּלֵיהֶם עֵץ וָאֶבֶן כֶּסֶף וְזָהָב אֲשֶׁר עִמָּהֶם׃}
}
{
	\eDtr1{
		And you’ve seen their disgraces and their idols, wood and
		stone, silver and gold, that were with them.
	}
}
{
}

\Verse{17}
{
	\hDtr1{פֶּן יֵשׁ בָּכֶם אִישׁ אוֹ אִשָּׁה אוֹ מִשְׁפָּחָה אוֹ שֵׁבֶט אֲשֶׁר לְבָבוֹ פֹנֶה הַיּוֹם מֵעִם יְהֹוָה אֱלֹהֵינוּ לָלֶכֶת לַעֲבֹד אֶת אֱלֹהֵי הַגּוֹיִם הָהֵם פֶּן יֵשׁ בָּכֶם שֹׁרֶשׁ פֹּרֶה רֹאשׁ וְלַעֲנָה׃}
}
{
	\eDtr1{
		In case there will be among you a man or woman or family
		or tribe whose heart is turning today from YHWH, our {\God},
		to go to serve those nations’ gods; in case there is among
		you a root bearing poison and wormwood,
	}
}
{
}

\Verse{18}
{
	\hDtr1{וְהָיָה בְּשׁמְעוֹ אֶת דִּבְרֵי הָאָלָה הַזֹּאת וְהִתְבָּרֵךְ בִּלְבָבוֹ לֵאמֹר שָׁלוֹם יִהְיֶה לִּי כִּי בִּשְׁרִרוּת לִבִּי אֵלֵךְ לְמַעַן סְפוֹת הָרָוָה אֶת הַצְּמֵאָה׃}
}
{
	\eDtr1{
		and it will be when he hears the words of this oath that
		he’ll feel himself blessed in his heart, saying, ‘I’ll
		have peace, though I’ll go on in my heart’s obstinacy,’ so
		as to annihilate the wet with the dry:
	}
}
{
%	\footnote{
%		oath. The word can mean either “oath” or “curse.” Here it can refer to
%		the covenant oath mentioned in v. 11 or to the covenant curses; or it
%		can mean the two together, since the covenant oath implicitly invokes
%		the curses on the one who takes it. Footnote 29:18. annihilate the wet
%		with the dry. Commentators and translators struggle with this phrase,
%		which may have been an idiomatic expression whose context is no longer
%		known. The answer may lie in the phenomenon that we have seen a number
%		of times whereby expressions from Genesis recur here in Deuteronomy. The
%		verb here (“annihilate” or “sweep away”) occurs four times in the matter
%		of Sodom and Gomorrah (Gen 18:23,24; 19:15,17) and in one other place in
%		the Torah, the destruction of Korah and his group (Num 16:26). In all
%		five occurrences there is an issue of whether good people will be
%		annihilated along with the bad. And they all involve fire coming from
%		God and destroying the offenders. The meaning of this expression here,
%		therefore, may be that as a strong fire burns up what is dry and
%		evaporates what is wet (or as a forest fire destroys both the dry and
%		fresh trees), so the obstinate person, or group in this case, may bring
%		down the community with them. This understanding is reinforced by the
%		fact that Sodom and Gomorrah are referred to just four verses later
%		(29:22).
%	}
}

\Verse{19}
{
	\hDtr1{לֹא יֹאבֶה יְהֹוָה סְלֹחַ לוֹ כִּי אָז יֶעְשַׁן אַף יְהֹוָה וְקִנְאָתוֹ בָּאִישׁ הַהוּא וְרָבְצָה בּוֹ כּל הָאָלָה הַכְּתוּבָה בַּסֵּפֶר הַזֶּה וּמָחָה יְהֹוָה אֶת שְׁמוֹ מִתַּחַת הַשָּׁמָיִם׃}
}
{
	\eDtr1{
		YHWH will not be willing to forgive him, because YHWH’s
		anger and His jealousy will then smoke against that man,
		and every curse that is written in this scroll will weigh
		on him, and YHWH will wipe out his name from under the
		skies.
	}
}
{
}

\Verse{20}
{
	\hDtr1{וְהִבְדִּילוֹ יְהֹוָה לְרָעָה מִכֹּל שִׁבְטֵי יִשְׂרָאֵל כְּכֹל אָלוֹת הַבְּרִית הַכְּתוּבָה בְּסֵפֶר הַתּוֹרָה הַזֶּה׃}
}
{
	\eDtr1{
		And YHWH will separate him from all of Israel’s tribes for
		bad, according to all the curses of the covenant that is
		written in this scroll of instruction.
	}
}
{
}

\Verse{21}
{
	\hDtr2{וְאָמַר הַדּוֹר הָאַחֲרוֹן בְּנֵיכֶם אֲשֶׁר יָקוּמוּ מֵאַחֲרֵיכֶם וְהַנּכְרִי אֲשֶׁר יָבֹא מֵאֶרֶץ רְחוֹקָה וְרָאוּ אֶת מַכּוֹת הָאָרֶץ הַהִוא וְאֶת תַּחֲלֻאֶיהָ אֲשֶׁר חִלָּה יְהֹוָה בָּהּ׃}
}
{
	\eDtr2{
		“And a later generation will say—your children who will
		come up after you, and the foreigner who will come from a
		far land, when they’ll see that land’s plagues and its
		illnesses that YHWH put in it,
	}
}
{
}

\Verse{22}
{
	\hDtr2{גּפְרִית וָמֶלַח שְׂרֵפָה כל אַרְצָהּ לֹא תִזָּרַע וְלֹא תַצְמִחַ וְלֹא יַעֲלֶה בָהּ כּל עֵשֶׂב כְּמַהְפֵּכַת סְדֹם וַעֲמֹרָה אַדְמָה (וצביים) [וּצְבוֹיִם] אֲשֶׁר הָפַךְ יְהֹוָה בְּאַפּוֹ וּבַחֲמָתוֹ׃}
}
{
	\eDtr2{
		brimstone and salt, all the land a burning, it won’t be
		seeded and won’t grow, and not any vegetation will come up
		in it, like the overturning of Sodom and Gomorrah, Admah
		and Zeboiim, which YHWH overturned in His anger and His
		fury—
	}
}
{
}

\Verse{23}
{
	\hDtr2{וְאָמְרוּ כּל הַגּוֹיִם עַל מֶה עָשָׂה יְהֹוָה כָּכָה לָאָרֶץ הַזֹּאת מֶה חֳרִי הָאַף הַגָּדוֹל הַזֶּה׃}
}
{
	\eDtr2{
		and all the nations will say, ‘For what did YHWH do
		something like this to this land? What is this big flaring
		of anger?’
	}
}
{
}

\Verse{24}
{
	\hDtr2{וְאָמְרוּ עַל אֲשֶׁר עָזְבוּ אֶת בְּרִית יְהֹוָה אֱלֹהֵי אֲבֹתָם אֲשֶׁר כָּרַת עִמָּם בְּהוֹצִיאוֹ אֹתָם מֵאֶרֶץ מִצְרָיִם׃}
}
{
	\eDtr2{
		“And they’ll say, ‘For the fact that they left the
		covenant of YHWH, their fathers’ {\God}, which He made with
		them when He brought them out from the land of Egypt.
	}
}
{
}

\Verse{25}
{
	\hDtr2{וַיֵּלְכוּ וַיַּעַבְדוּ אֱלֹהִים אֲחֵרִים וַיִּשְׁתַּחֲווּ לָהֶם אֱלֹהִים אֲשֶׁר לֹא יְדָעוּם וְלֹא חָלַק לָהֶם׃}
}
{
	\eDtr2{
		And they went and served other gods and bowed to them,
		gods whom they hadn’t known and He hadn’t allocated to
		them.
	}
}
{
}

\Verse{26}
{
	\hDtr2{וַיִּחַר אַף יְהֹוָה בָּאָרֶץ הַהִוא לְהָבִיא עָלֶיהָ אֶת כּל הַקְּלָלָה הַכְּתוּבָה בַּסֵּפֶר הַזֶּה׃}
}
{
	\eDtr2{
		And YHWH’s anger flared at that land, to bring over it
		every curse that was written in this scroll.
	}
}
{
}

\Verse{27}
{
	\hDtr2{וַיִּתְּשֵׁם יְהֹוָה מֵעַל אַדְמָתָם בְּאַף וּבְחֵמָה וּבְקֶצֶף גָּדוֹל וַיַּשְׁלִכֵם אֶל אֶרֶץ אַחֶרֶת כַּיּוֹם הַזֶּה׃}
}
{
	\eDtr2{
		And YHWH plucked them from their land in anger and in fury
		and in great rage, and He threw them into another land, as
		it is this day.’
	}
}
{
}

\Verse{28}
{
	\hDtr1{הַנִּסְתָּרֹת לַיהֹוָה אֱלֹהֵינוּ וְהַנִּגְלֹת לָנוּ וּלְבָנֵינוּ עַד עוֹלָם לַעֲשׂוֹת אֶת כּל דִּבְרֵי הַתּוֹרָה הַזֹּאת׃}
}
{
	\eDtr1{
		“The hidden things belong to YHWH, our {\God}, and the
		revealed things belong to us and to our children forever,
		to do all the words of this instruction.
	}
}
{
%	\footnote{
%		The hidden things belong to YHWH. Moses tells the people that there are
%		things that are solely in the divine realm, hidden from humans. Humans’
%		task is here in this world: to live a certain kind of life. He breaks
%		off from this matter and turns to a discussion of future exile and
%		return to God (30:1–10). Then he returns to this picture of what humans
%		must do, and how it is not something enigmatic, in v. 11.
%	}
}
